\chapter{Related Work and Discussion}
\label{sec:discussion}

Chapter~\ref{sec:evaluation} established what this thesis's four parsers
do and how they perform, against each other and against
\texttt{regex}/RE2. This chapter steps back from those numbers to ask
three different questions: where does this work sit relative to other
approaches to the same problem (Section~\ref{sec:related-work}); what did
translating an academic Haskell reference into a systems language
actually involve (Section~\ref{sec:rust-vs-haskell}); and what did
building the tracing-based diagnostics facility that Chapter~\ref{sec:conclusion}
frames as a first step toward a fuller system actually reveal about
derivative-based parsing's own structure (Section~\ref{sec:mismatch-analysis})?
A final section (Section~\ref{sec:optimization-challenges}) is
retrospective rather than forward-looking, unlike
Section~\ref{sec:further-performance-optimizations}: it is about what was
actually tried during this thesis's development, including two attempts
that looked like improvements and were not, not about what remains to be
tried.



\section{Related Work}
\label{sec:related-work}

\subsection{POSIX-Compliant Engines}
\label{sec:posix-engines-related-work}

Two established engines reach POSIX leftmost-longest correctness by a
mechanism entirely different from this thesis's derivative-based one.
Laurikari's tagged NFA construction \citep{laurikari2000} propagates
\emph{tags} -- markers recording where a submatch begins or ends -- through
NFA transitions, then determinizes the tagged NFA into a tagged DFA
(TDFA) that can be compiled once and reused, with submatch information
falling out of which tagged transitions a run actually took. Kuklewicz's
\texttt{regex-tdfa} \citep{kuklewicz_tdfa} implements this construction in
Haskell, offering POSIX leftmost-longest matching with the same
compile-once-run-many-times shape Section~\ref{sec:performance-comparison}
attributes to \texttt{regex}/RE2, not to any parser in this thesis.

The contrast worth drawing is not who is faster -- Chapter~\ref{sec:evaluation}
already establishes that this thesis's parsers, walking live tree/bit
structures on every match, cannot approach a compiled automaton's speed,
tagged or not -- but \emph{where} each approach puts the information a
POSIX parse tree needs. A tagged DFA carries it as annotations on
automaton transitions, decided once during compilation and replayed on
every subsequent match. This thesis's constructions carry it either in an
explicit backward pass over the derivative sequence
(\texttt{deriv\_std}, Section~\ref{sec:core-algorithm}) or in bits
threaded through the derivative computation itself
(\texttt{deriv\_bc}, Section~\ref{sec:bitcoded-optimization-theory}) --
recomputed on every match, never cached. That TDFA-based engines
demonstrate POSIX correctness and compiled, reusable automata are not
mutually exclusive is precisely why
Section~\ref{sec:derivative-automaton-caching} treats caching the
derivative automaton across matches as a real, precedented direction
rather than a speculative one -- tagged NFAs are an existence proof that
\emph{some} way of attaching submatch information to a compiled,
cacheable automaton is possible, even though this thesis does not
establish that a derivative-based version of that attachment is.

\subsection{RE2 and Rust's \texttt{regex} Crate}
\label{sec:automaton-engines-related-work}

Section~\ref{sec:compared-engines} and Section~\ref{sec:performance-comparison}
already compare this thesis's parsers against \texttt{regex} and RE2
empirically, on real numbers. What is worth adding here is the
conceptual placement those numbers do not by themselves argue for: three
different families -- automaton \emph{simulation}
(\texttt{regex}/RE2, Section~\ref{sec:rust-regex-crate-eval}-\ref{sec:re2-eval}),
automaton \emph{simulation with tags}
(Section~\ref{sec:posix-engines-related-work}), and algebraic
\emph{transformation} (this thesis, and
Section~\ref{sec:other-derivative-work} below) -- are three different
answers to the same underlying tension between correctness/expressiveness
and compiled-time performance, not three points on a single ranking.
Cox's own account of why backtracking engines are avoided in favor of
automaton simulation \citep{cox2007} is the standard reference for why
the first family made the architectural choice it did; this thesis's
Chapter~\ref{sec:motivation} already leans on it for the same reason.
Positioning derivative-based parsing as a third family, rather than as an
underperforming instance of the first, is the more accurate reading of
where Chapter~\ref{sec:evaluation}'s gap actually comes from.

\subsection{Other Derivative-Based Approaches}
\label{sec:other-derivative-work}

Owens, Reppy, and Turon revisit Brzozowski's original derivative
construction directly \citep{owens2009}, already introduced in
Section~\ref{sec:derivative-automaton-caching} as the literature this
thesis's own "cache the automaton" direction would build on: derivatives
used to \emph{compile} a DFA, each distinct derivative (up to a
normalizing equivalence) a state, each character a transition, producing
smaller automata than the classical McNaughton--Yamada construction and
extending cleanly to large character sets. This is the plain-derivative,
membership-only case; Section~\ref{sec:reference-scope} below notes that
Sulzmann's own teaching-page material \citep{sulzmann_slides} sketches
the closer, partial-derivative analogue this thesis's own
\texttt{pderiv\_bc} could in principle build on -- a construction closer
to hand than Owens, Reppy, and Turon's alone suggested when
Section~\ref{sec:derivative-automaton-caching} was first written.

Might, Darais, and Spiewak extend the same foundational idea in a
different direction, from regular expressions to context-free grammars
\citep{mightdaraisspiewak2011}: derivatives of a CFG, computed lazily,
memoized, and tied together with fixed points to keep an otherwise
circular definition well-founded. This is a more ambitious generalization
than the one this thesis undertakes -- Sulzmann and Lu's own extension
\citep{sulzmannlu2014} stays within regular languages, adding POSIX
disambiguation rather than generalizing the grammar class -- but it
establishes that "derivative" as a parsing-construction technique
generalizes well beyond the regular case this thesis, and the original
Brzozowski construction it builds on, are both scoped to.

\subsection{Kuklewicz's Observation on Ordering-Rule Ambiguity}
\label{sec:kuklewicz-observation}

Kuklewicz's practical experience building \texttt{regex-tdfa}
(Section~\ref{sec:posix-engines-related-work}) is widely cited, in
secondary accounts of the POSIX-matching literature, for a specific,
uncomfortable finding: engines that each describe themselves as
"POSIX-compliant" do not, in practice, agree with each other on
genuinely ambiguous input. \citet{sulzmannlu2014} state this directly in
their own related-work discussion, citing \citet{kuklewicz_regex_posix}
by name: \emph{"Kuklewicz observes that almost all POSIX implementations
are buggy which is confirmed by our own experiments."} POSIX's own
textual specification of leftmost-longest matching -- the same
recursive, subexpression-by-subexpression ordering rule this thesis's
own Definition 1 (Section~\ref{sec:posix-longest-leftmost}) formalizes
as rules (A1)--(A4)/(K1)--(K4) -- leaves enough room, read informally,
for different implementers to make different, individually
self-consistent, mutually inconsistent choices about how ambiguity
nested inside ambiguity should resolve; Kuklewicz's own wiki page
catalogues exactly this kind of disagreement across the POSIX-labeled
engines available to him at the time, not a single bug in a single
implementation.

What this thesis's own construction offers in exactly this situation is
a \emph{formal} answer to the ordering question rather than an
operational one: Definition 1 (Section~\ref{sec:posix-longest-leftmost}) is a
mathematical relation on parse trees, proved (Theorem 1,
Section~\ref{sec:correctness-properties}) to be exactly what
\texttt{deriv\_std}/\texttt{deriv\_bc} compute -- not merely what one
particular codebase's control flow happens to produce. Two engines
implementing the same formal ordering relation correctly cannot disagree
with each other on an ambiguous input in the way Kuklewicz's practical
experience suggests differently-engineered "POSIX" engines can.



\section{Rust versus Haskell}
\label{sec:rust-vs-haskell}

\subsection{Language Characteristics of the Reference Implementation}
\label{sec:haskell-language-characteristics}

The reference material this thesis translates from -- both
flops14-extended's own presentation \citep{sulzmannlu2014} and
Sulzmann's fuller teaching-page Haskell transcript \citep{sulzmann_slides}
-- leans on several Haskell features with no direct Rust counterpart,
each already surfaced individually at the point in this thesis where it
first mattered; this section collects them.

\paragraph{Algebraic data types with blanket \texttt{deriving}.} The
teaching page's own regex type,
\begin{verbatim}
data R = Phi | Eps | L Char | Seq R R | Choice R R | Star R
  deriving (Show, Eq, Ord)
\end{verbatim}
derives structural equality, ordering, \emph{and} printing in one clause,
whether or not a given piece of code actually uses all three. Section~\ref{sec:regex-rust-implementation}'s
\texttt{Regex} derives only \texttt{Debug}, \texttt{Clone},
\texttt{PartialEq}, \texttt{Eq}, and \texttt{Hash} -- no \texttt{Ord} --
because nothing in this thesis's Rust port ever compares two
\texttt{Regex} values for ordering; Rust's per-trait, opt-in
\texttt{derive} makes that omission a visible, deliberate fact about the
codebase, where the reference's blanket \texttt{deriving} leaves it
implicit.

\paragraph{GADT syntax used without a type index.} The reference's parse
tree type is written as a GADT,
\begin{verbatim}
data U where
  Nil :: U
  Empty :: U
  Letter :: Char -> U
  LeftU :: U -> U
  RightU :: U -> U
  Pair :: (U,U) -> U
  List :: [U] -> U
\end{verbatim}
but no constructor's return type is ever refined to anything other than
plain \texttt{U} -- this is an ordinary sum type written in GADT syntax,
not a use of GADTs' actual type-level power. Section~\ref{sec:parsetree-and-flatten}'s
\texttt{ParseTree} enum needs nothing beyond an ordinary Rust
\texttt{enum} to match it.

\paragraph{List comprehensions.} Both \texttt{pDeriv} and \texttt{pDerivBC}
are written with list comprehensions throughout, e.g.\
\texttt{[ Seq r' s | r' <- pDeriv x r ]}. Every one of this thesis's Rust
equivalents (\texttt{pderiv\_bc}, Section~\ref{sec:pderiv-bc-theory-rust};
the plain \texttt{pderiv} matcher, Section~\ref{sec:matcher-benchmark-fairness})
writes the identical transformation as an iterator chain --
\texttt{.into\_iter().map(...).collect()} -- a direct, mechanical
translation with no semantic gap, differing only in surface syntax.

\paragraph{\texttt{Data.List.nub}/\texttt{nubBy}.} The reference reaches
for generic, standard-library list deduplication casually --
\texttt{nub} in \texttt{pDeriv}, \texttt{nub2 = nubBy (\textbackslash(r,\_) (s,\_) -> r == s)}
in \texttt{pDerivBC}. Rust's standard library has no equivalent; every
dedup in this thesis's Rust port (\texttt{nub2},
Section~\ref{sec:pderiv-bc-theory-rust}; the \texttt{match\_pderiv}
investigation, Section~\ref{sec:performance-comparison}) is a
hand-written choice between \texttt{HashSet} and \texttt{Vec}, with the
performance trade-offs that choice carries. Worth noting explicitly:
GHC's own \texttt{nub}/\texttt{nubBy} are themselves the na\"ive,
quadratic list-based implementation, not a hashed one -- the reference's
brevity here does not imply the reference is fast at this operation, only
that its own text never has to engage with the question at all, where
every Rust translation of it does.

\paragraph{Laziness.} The teaching page's simplification function,
\begin{verbatim}
simp r = let s = simp2 r in if s == r then r else simp s
\end{verbatim}
recomputes \texttt{simp2} and re-checks equality in an explicit,
unbounded recursive loop with no argument, in the source itself, for why
this terminates in practice beyond the informal claim that fixpoints are
reached quickly. Laziness makes writing this kind of "keep going until
nothing changes" recursion low-risk in Haskell in a way it is not in
Rust's eager, natively-stacked evaluation -- the same distinction
Section~\ref{sec:stack-depth}'s \texttt{parse\_recursive}/\texttt{parse\_loop}
split, and \texttt{examples/demo\_crash.rs}'s existence, exist to manage
directly rather than leave implicit.

\paragraph{\texttt{error} for impossible cases.} The reference's partial
functions signal "this case cannot arise in well-typed use" with
\texttt{error}, e.g.\ \texttt{mkEpsBC (L l) = error "mkEpsBC, letter impossible"}.
This thesis's Rust equivalents use \texttt{panic!} for the identical
purpose and, in several cases, the identical case -- \texttt{mk\_eps\_bc.rs}'s
panics on \texttt{Phi}/\texttt{Lit} mirror this exact reference line, not
merely the general pattern.

\subsection{What the Reference Provides vs. What It Leaves Open}
\label{sec:reference-scope}

flops14-extended \citep{sulzmannlu2014} provides the core algorithm
(\texttt{mkEps}/\texttt{inj}/\texttt{parse}, Figure 3), the bit-coded
refinement (Figures 4--6), a Haskell reference implementation, a
POSIX-correctness proof (Theorem 1, Section~\ref{sec:correctness-properties}), and a
linear-time complexity claim -- but explicitly leaves space complexity
and a full proof of it as future work
(Section~\ref{sec:bitcoded-optimization-theory}), a gap
Section~\ref{sec:pathological-inputs} shows this thesis's own
\texttt{deriv\_bc} inherits concretely, not merely by citation, on
$(a+aa)^{*}$.

What the paper does not provide at all: any partial-derivative-based
POSIX construction (motivating Chapter~\ref{sec:partial-derivative-parser}'s
own synthesis, since neither primary source defines one --
Section~\ref{sec:partial-derivative-parser-theory}); a benchmarked
implementation in a systems language; any tracing or diagnostics
facility (Section~\ref{sec:mismatch-analysis} below); and any discussion
of substring-search or padding semantics at all --
\texttt{docs/SUBSTRING\_SEARCH.md}'s account of the
$\Sigma^{*} \cdot r \cdot \Sigma^{*}$ padding problem was resolved
through direct email correspondence with Sulzmann and Lu themselves,
precisely because the paper is silent on it and the question turned out,
by their own account in that correspondence, to still be open.

Two findings surfaced only while researching \emph{this} chapter are
worth recording here rather than left implicit in earlier chapters that
did not have occasion to mention them.

First, Sulzmann's teaching-page material provides \emph{more} than the
paper alone: a working prototype of a bit-coded partial-derivative NFA
construction, \texttt{buildBC\_NFA}/\texttt{runBC\_NFA}, that caches
every reachable \texttt{(residual, bits)} transition once and replays it
across matches -- exactly the "cache the automaton" direction
Section~\ref{sec:derivative-automaton-caching} discusses as a
literature-grounded but unattempted direction for this thesis, except
specific to the \emph{partial}-derivative, bit-coded case this thesis's
own \texttt{pderiv\_bc} most directly extends, rather than the
plain-derivative case Owens, Reppy, and Turon treat. This thesis does not
port or benchmark this prototype; it is recorded here as a closer
starting point than was known to exist when
Section~\ref{sec:derivative-automaton-caching} was first written, not as
work this thesis completes.

Second, the teaching page documents its own bug and fix in place, as a
comment left directly in the source:
\begin{verbatim}
pDerivBC x (Cat r s)
  | nullable r = nub2 $ [ (Cat r' s,bs) | (r',bs) <- pDerivBC x r ]
                        ++
                        -- (pDerivBC x s) -- bug spotted by clause
                        [ (s', mkEpsBC r ++ bs) | (s',bs) <- pDerivBC x s ]
\end{verbatim}
The commented-out line omits prefixing the second strand's bits with
\texttt{mkEpsBC r} -- exactly the incremental bit-fusing step
Section~\ref{sec:partial-derivative-parser-bitcoded-theory} transcribes
as part of \texttt{pDerivBC}'s own equations. This thesis's
\texttt{pderiv\_bc.rs} (Section~\ref{sec:pderiv-bc-theory-rust}) was
checked directly against this source line while writing this chapter and
matches the corrected version, not the commented-out one. The concrete
point is not that a bug was avoided -- it was, but that alone would be
unremarkable -- it is that the reference material translated throughout
this thesis was itself an evolving, imperfect artifact with its own
visible debugging history, not a finished, unimpeachable ground truth;
faithfully translating it required reading closely enough to tell an
author's corrected intent from a mistake still sitting, commented out, in
the same block.

\subsection{Translation Strategy}
\label{sec:translation-strategy}

Five axes recur across every individual translation decision this thesis
made, each already introduced locally at the point it first mattered;
collecting them here is what the forward-reference from
Section~\ref{sec:matching-vs-parsing} promised.

\textbf{Evaluation strategy.} Haskell's laziness sidesteps concerns Rust's
eager evaluation forces into the open: unbounded recursion depth
(Section~\ref{sec:stack-depth}'s \texttt{parse\_recursive}/\texttt{parse\_loop}
split) and unbounded repetition counts (the frontend's \texttt{MAX\_BOUND}
repetition cap, needed because a bound like \texttt{\{1586,\}} would
otherwise unroll into an impractically large \texttt{Regex} eagerly,
where Haskell's laziness could in principle defer that cost
indefinitely).

\textbf{Control flow.} Haskell's pattern-matching equations and Rust's
\texttt{match} arms are close to a direct, one-to-one translation for
case structure itself -- every derivative, nullability, \texttt{mkEps},
and \texttt{decode} function in this thesis follows its reference's case
structure exactly (Section~\ref{sec:regex-rust-implementation}). The
friction is elsewhere on this list, not here.

\textbf{Data types.} The recursive ADT needs \texttt{Box}
(Section~\ref{sec:regex-rust-implementation}); the GADT-syntax parse tree
needs nothing beyond a plain \texttt{enum}
(Section~\ref{sec:haskell-language-characteristics} above). One further
instance, already noted where it first arose
(Section~\ref{sec:parsetree-and-flatten}): the reference's \texttt{Nil}
constructor exists because \texttt{mkEps}, as a total Haskell function,
still needs \emph{some} value for the case its own type does not rule
out by construction (\texttt{mkEps Phi = Nil}, since nothing prevents
\texttt{mkEps} being applied to a non-nullable \texttt{Phi} at the type
level, even though it never legitimately is). \texttt{Option<ParseTree>::None}
represents the analogous "no parse tree" case at this thesis's own
top level directly, rather than needing a dedicated sentinel constructor
threaded through every intermediate function's return type.

\textbf{Pattern matching.} Guards translate directly
(\texttt{| x == y = ... | otherwise = ...} to Rust's own guard syntax).
List patterns over the input word (\texttt{(x:xs)}) become iterator-based
consumption (\texttt{input.chars()}) -- semantically identical, since
both decompose a sequence one element at a time, but idiomatically
different at every call site.

\textbf{Abstraction.} Higher-order list-comprehension lambdas become Rust
closures passed to iterator adaptors -- functionally equivalent,
syntactically heavier, and, per
Section~\ref{sec:haskell-language-characteristics}'s \texttt{nub}
discussion, occasionally hiding a real implementation decision (which
data structure backs a deduplication step) behind what reads, in the
reference, as one uniform standard-library call.

\textbf{Memory model.} Garbage collection and implicit structural sharing
in Haskell against explicit ownership, \texttt{Box}, and \texttt{Clone}
in Rust is the axis every other item on this list ultimately reduces to.
It is also the one Section~\ref{sec:further-performance-optimizations}'s
\texttt{Rc}-in-place-of-\texttt{Box} proposal and
\texttt{pderiv\_std}'s injection-closure design
(\texttt{Rc<dyn Fn(ParseTree) -> ParseTree>},
Section~\ref{sec:pderiv-std-theory-rust}) both exist specifically to
manage -- the latter is, concretely, Rust's explicit stand-in for a
closure Haskell's laziness would otherwise capture and defer for free.

\subsection{Translation Insights}
\label{sec:translation-insights}

What carried over easily: the \emph{case structure} of every function
defined by structural recursion on the regex ADT. This is not a surprise
given Section~\ref{sec:regex-rust-implementation}'s observation that
both languages' pattern matching enforces the same "every constructor
handled" totality guarantee -- a derivative or nullability function
defined case-by-case in Haskell has essentially one correct Rust shape to
translate to, and every function in this thesis following that shape
(\texttt{deriv}, \texttt{nullable}, \texttt{mk\_eps}, \texttt{mk\_eps\_bc},
\texttt{decode}, \texttt{pderiv\_bc}) bears this out directly.

What did not: anything touching sharing or laziness, without exception.
The \texttt{parse\_recursive}/\texttt{parse\_loop} split
(Section~\ref{sec:stack-depth}) has no Haskell-reference analogue at all,
since the reference's own recursion depth concerns, if any, are GHC's to
manage, not a translation decision this thesis had to make explicitly.
The \texttt{pderiv\_std} injection-closure design exists purely to
recreate, explicitly, capture-and-defer behavior Haskell's laziness
provides without any corresponding design decision visible in the
reference text at all. And every \texttt{nub}/deduplication call site
needed a deliberate choice of backing data structure that the reference's
own \texttt{nub}/\texttt{nubBy} calls never surface as a decision --
Section~\ref{sec:performance-comparison}'s \texttt{match\_pderiv} finding
(a redundant, needlessly repeated \texttt{HashSet} costing roughly a
fifth of that benchmark's total time) is a direct instance of a cost the
reference's brevity had made invisible, discovered only once the
translation had to make the choice concrete and pay for it.

Finally, Section~\ref{sec:reference-scope}'s bug-fix comment is itself a
translation insight, not just a curiosity: faithfully translating
reference code is not a mechanical process of transcribing equations,
because the reference material itself carries its own editing history,
including corrections its author made and left visible rather than
silently overwrote. Reading closely enough to translate the corrected
\texttt{pDerivBC} rather than the commented-out mistake sitting beside it
was a real, if small, act of judgment this thesis's own \texttt{pderiv\_bc.rs}
had to get right, not something guaranteed by translating "the reference"
as an undifferentiated whole.



\section{Mismatch Analysis via Derivative Tracing}
\label{sec:mismatch-analysis}

\subsection{The Gap}
\label{sec:diagnostics-gap}

A production regex engine's answer to a failed match is, ordinarily, a
boolean: no. \texttt{regex}/RE2's \texttt{is\_match} returning
\texttt{false} carries no information about \emph{why} -- which part of
the pattern first ran out of options, at which position in the input, or
what the closest successful alternative would have needed to look like.
This is not a design flaw in those engines; it is a direct consequence of
Section~\ref{sec:performance-comparison}'s own architectural point --
once a pattern is compiled to a table-driven automaton, the
"reasoning" that produced that table has already been discarded by the
time a match is attempted, and there is nothing left, at match time, to
report on.

\texttt{src/diagnostics/replay.rs}'s \texttt{error\_report}/\texttt{find\_failure}
give a concrete sense of what is missing: given a failed match, they walk
the recorded derivative sequence to find exactly the position where no
further derivative step could succeed, and render a caret pointing at it
directly against the input string -- the kind of answer "no" alone
cannot give, produced from data an automaton-based engine has no
equivalent of retaining at all.

\subsection{Deriving Diagnostics for Free from the Algorithm Itself}
\label{sec:diagnostics-for-free}

The point worth making precisely, not just informally: derivative-based
parsing already computes an explicit sequence
$r_0, r_1, \ldots, r_n$ (Section~\ref{sec:core-algorithm}) as its
\emph{ordinary} operation, whether or not anything downstream ever
inspects it. Recording that sequence for diagnostic purposes costs
nothing algorithmically beyond the recording itself -- no separate
instrumentation pass, no auxiliary structure the matching process would
not otherwise build. An automaton-based engine has no equivalent
sequence to record at match time at all: whatever "reasoning" a Thompson
NFA or lazy DFA construction performed happened once, at compile time,
and is baked into the transition table a later match merely walks: the
information a diagnostic would need has already been compiled away by
the time a specific match is attempted, and reconstructing anything like
it needs separately-built tracing infrastructure bolted on afterward, not
already sitting in the algorithm's own intermediate state the way it
does here.

This is, concretely, one of this thesis's own contributions
(Section~\ref{sec:contributions}) rather than an incidental convenience:
the tracing facility \texttt{src/diagnostics/} builds on is a direct
consequence of choosing a derivative-based construction in the first
place, not a feature that would have been equally easy to add to an
automaton-based engine as an afterthought.

\subsection{Design}
\label{sec:diagnostics-design}

\texttt{DiagLevel} (\texttt{src/diagnostics/mod.rs}) has four values --
\texttt{Off} (0), \texttt{Basic} (1, regex/input/match/tree/error caret),
\texttt{Verbose} (2, \texttt{Basic} plus timing, expression count, and
construction steps), and \texttt{Debug} (3, the full derivation trace,
written to a file if \texttt{REGEX\_DIAG\_REPORT} is set) -- selected via
the CLI's \texttt{--diag} flag or the \texttt{REGEX\_DIAG} environment
variable, and applies only to \texttt{parse}, not \texttt{match}: a
Boolean matcher has no parse tree, and hence nothing at levels 1--3 to
report beyond what the trace itself would show, which level 3 already
provides directly. Each of the four parser variants renders its own
trace independently rather than sharing one generic renderer --
\texttt{diag\_2/\{deriv\_std,deriv\_bc,pderiv\_std,pderiv\_bc\}.rs} and
the analogous four files under \texttt{diag\_3/} -- since what is worth
showing at each level differs by construction (a bit-coded parser's trace
is bit-strings and \texttt{ARegex} states; the tree-injection parsers'
traces are \texttt{Regex} derivative sequences and, at level 3, the
backward injection steps that reconstruct a tree from them).
\texttt{docs/DIAG.md} documents concrete output at every level, for every
parser, on both successful and failing matches.

\subsection{Scope and Limits}
\label{sec:diagnostics-scope-limits}

What the current system does not do, stated plainly rather than left to
be discovered: it reports where a match failed, not why in any
higher-level sense -- no ranked list of near-miss alternatives, no
natural-language explanation of what the input would have needed to
contain, no interactive step-through letting a user replay the trace one
decision at a time rather than reading it as a fixed report.
Section~\ref{sec:mismatch-diagnostics-future} (Chapter~\ref{sec:conclusion})
takes up what a fuller system addressing these gaps would need to add,
building on the tracing facility this section describes rather than
replacing it.



\section{Optimization Challenges and Insights}
\label{sec:optimization-challenges}

Section~\ref{sec:further-performance-optimizations} is forward-looking:
concrete, scoped work this thesis identifies but does not attempt. This
section is the opposite -- a retrospective account of optimization work
actually undertaken during this thesis's development, including two
attempts that looked like real improvements and were rejected once tested
against adversarial input, not just the one that shipped. The throughline
across all three case studies below is the same: "faster on the
benchmark in front of you" and "faster, full stop" are not the same
claim, and in every case here the gap between them was found only by
testing directly against an adversarial input, not by reasoning about
asymptotic complexity alone.

\paragraph{The $(a+aa)^{*}$ blowup on \texttt{deriv\_bc}.} First
encountered as an unexplained slowdown while benchmarking
\texttt{ambiguous\_star\_a\_or\_aa} (below), not predicted in advance:
\texttt{deriv\_bc}'s \texttt{simp} (Section~\ref{sec:simplification-rules-theory})
turned out to be the \emph{slowest} of four parsers on this pattern,
despite being the only one of the four with an explicit simplification
pass at all. \texttt{docs/DERIV\_BC.md} documents the full investigation:
\texttt{simp}'s deduplication is faithful to Figure 6's own \texttt{nub}
-- bit-inclusive, not shape-only -- and POSIX correctness genuinely
requires this, unlike \texttt{pderiv\_bc}'s shape-only \texttt{nub2}
(Section~\ref{sec:partial-derivative-parser-theory}), because two
branches with identical residual shape but different accumulated bits
are different candidate parse histories, and discarding one by shape
alone risks discarding the POSIX-correct one. The consequence, verified
directly rather than assumed: on $(a+aa)^{*}$, whose ambiguity does not
collapse under this dedup, the annotated tree's node count grows as
exact Fibonacci numbers, and the resulting right-associative \texttt{Alt}
chain overflows the stack around $n=20$--$22$, well before wall-clock
cost alone would have made the pattern impractical.

\paragraph{The \texttt{ambiguous\_star\_a\_or\_aa} benchmark fix.} A
direct consequence of the finding above, and a case study in fixing the
right thing: \texttt{benches/bench\_posix.rs}'s size range for this group
originally jumped straight from 10 to 50, past every one of the four
parsers' actual safety margin -- not just \texttt{deriv\_bc}'s stack
overflow, but \texttt{recursive}/\texttt{pderiv\_bc}/\texttt{pderiv\_std}
all failing to return in reasonable wall-clock time on the same range.
The fix was to the benchmark's own parameters (Section~\ref{sec:ambiguous-star-results}
reports the resulting, safe numbers), not to the parsers -- the blowup
above is a real, inherent property of this pattern under
\texttt{deriv\_bc}'s construction, not a bug a code change could remove
without weakening the POSIX guarantee \texttt{simp}'s dedup exists to
preserve.

\paragraph{The \texttt{match\_pderiv} \texttt{HashSet} fix.} The most
instructive of the three, because two faster-looking alternatives were
built and rejected before the change that actually shipped
(Section~\ref{sec:performance-comparison}). \texttt{match\_pderiv}
originally deduplicated via a live \texttt{HashSet<Regex>} at every
\texttt{Alt}/\texttt{Seq}/\texttt{Star} node during its own recursive
derivative computation, on top of the one \texttt{HashSet} its
top-level, per-character loop already needed for correctness -- the
inner deduplication was strictly redundant with, and subsumed by, the
outer one. The first fix attempted, dropping deduplication inside
\texttt{pderiv} entirely, looked correct by every test in the existing
suite and was roughly $57\%$ faster on the real corpus benchmark -- and
was rejected only because it was tested directly against $(a+aa)^{*}$,
where it reproduced the exact Fibonacci blowup described above (over two
million states by 30 characters). A second attempt, deduplicating once
per node locally rather than not at all -- mirroring \texttt{pderiv\_bc}'s
own \texttt{nub2} -- looked like a reasonable middle ground and was
tested the same way: it very nearly reproduced the blowup too, since
local deduplication within one \texttt{pderiv} call does not bound the
union across the many \emph{different} top-level states the outer loop
tracks simultaneously, which is exactly the property the original,
unmodified outer \texttt{HashSet} was providing. The version actually
shipped removes only the demonstrably redundant inner \texttt{HashSet}
and nothing else, verified safe on $(a+aa)^{*}$ up to 200 characters
(identical to the pre-change implementation) and a real, smaller,
honest $21\%$ faster on the corpus benchmark -- smaller than either
rejected alternative's number, because it is the only one of the three
that does not also remove a property those alternatives never noticed
they were removing.

The lesson these three case studies share is not "always benchmark" in
the abstract -- every parser in this thesis was already benchmarked
throughout its development -- but specifically that a benchmark
measuring one workload (a real-world corpus sample, in
\texttt{match\_pderiv}'s case; a modest size range, in
\texttt{ambiguous\_star\_a\_or\_aa}'s) cannot certify safety on a
workload it does not include. Every real fix in this section, and every
rejected alternative, was decided by constructing the specific
adversarial input each construction's own theory predicted would be
worst-case, not by the existing test suite passing or a synthetic
benchmark improving.
